\section{导言：问题的坐标与本文的任务}

2026 年 10 月 6 日，OpenAI 在一个公开仓库中放出了内部前沿模型生成的数学手稿合集，共 372 个结果族；翌日的修订记录撤回了其中 3 篇并修订 14 篇，其余保留，此后按版本号管理修正\upcite{openairelease2026,openairepo2026}。这批材料里，第 074 族的体量不算最大，位置却最扎眼：它只有两篇稿件，主题是挂谷（Kakeya）猜想。一篇 97 页，声称证明三维挂谷极大函数猜想；另一篇 175 页，声称证明四维挂谷集的 Hausdorff 维数等于 $4$\upcite{openai3d2026,openai4d2026}。作者栏里没有个人署名，只有一个机构名。

值得为这一族单独写一份综述，理由不只是“AI 又解了一个大问题”。挂谷猜想的三维情形在 2025 年 2 月刚被王虹与 Zahl 解决，2026 年 7 月费城国际数学家大会上，这项工作成为王虹获菲尔兹奖的主要成果之一\upcite{wangz2025,xinhua2026,pku2026}。这是一条仍在高速推进、社区记忆新鲜的研究线，OpenAI 手稿直接踩在这条线的最前沿，而不是挑选一个冷门角落。另一方面，两篇稿件都没有进入仓库的 Lean 形式化清单，也没有经过同行评审；仓库自己的措辞是“部分未形式化的结果可能存在问题”\upcite{openairepo2026}。结果的真伪悬置，方法的面貌却已经摊开。这正是综述可以介入的时机。

本文的任务有六项。其一，把发布事实、稿件参数、依赖链条与验证状态整理成一份可核对的档案（第 4 章），其中最关键的一条是：四维稿件的核心输入取自三维姊妹篇的一条加权板条引理，两篇稿件在验证意义上是一体的，不能拆开评估。其二，逐层重建四维证明的技术骨架（第 5 章），覆盖从任意挂谷集到加权模型的归约、“地平线”机制、图卡体系、极值临界路径与三窄度区间的分类讨论。其三，做技术谱系分析（第 6、7 章）：把全部核心技术按直系继承、门内改造、跨门派借入、谱系内首次系统使用四类逐一归档，并给出证据与置信度。其四，把这条研究线上的主要方法放进一张统一比较矩阵（第 8 章）。其五，分析验证缺口（第 9 章），说明“未经形式化、思想未经独立推演”对这份工作的具体含义。其六，给出核验路线图与下一步方向建议（第 10 章），并在收束处给出修正后的判断（第 11 章）。

本文不做三件事。不做 175 页的逐行核验——那需要一支熟悉四维管族几何的验证小组，超出单篇综述的职责与能力；不复述稿件中的全部定义与引理编号，只保留理解机制所必需的骨架；不代替数学界裁决正确性，所有涉及稿件的结论一律写作“手稿声称”。

材料与方法是本文可信度的全部来源，因此一并说明。一手材料包括两篇稿件 PDF 全文、仓库的目录清单、Lean 形式化清单与修订史、OpenAI 官方发布说明\upcite{openai3d2026,openai4d2026,openairepo2026,openairelease2026}；历史与工具箱基线取自 Zahl 的挂谷综述（覆盖 2000--2025 年）\upcite{zahl2025survey}；为判断“哪些工具是门内已用、哪些是外门派借入”，本文对 Wang--Zahl 的三维论文与 Guth--Wang--Zahl 的简化证明两份门内主力文献做了全文关键词扫描，结果在第 7 章给出。事实性断言以上标编号标注出处；属于本文分析判断的，在文中明示，并按〔置信度：高/中/低〕标注。技术性转述以稿件自身的章节与命题编号为准（如 Theorem 2.17、Proposition 6.8），读者可以直接按编号定位原文核对，这也用于防止转述层引入不必要的失真。

导言的最后一段话说明本文的立场。挂谷猜想在四维是否已经被解决，本文不给出结论；本文给出的是这份稿件的结构、来源、缺口与可能去处。一份正确但无人理解的证明，和一份无人能验证的声称，在数学共同体里的地位完全不同；把二者区分开，是接下来十章要做的事。